% =====================================================================
%  Banco complementar --- Cargas Eletricas  (REVISADO)
%  Substitui a versao gerada por template (2 clones em N3 e 8 questoes
%  conceituais marcadas como N4).
%  RESTRICAO SEVERA: o cap01 e extenso (31 questoes) e ja cobre
%  praticamente todos os temas das N4 antigas: contato entre esferas
%  identicas (varias vezes), contatos sucessivos, esferas de raios
%  diferentes, eletroscopio, inducao com aterramento, verificacao de
%  quantizacao (Q = 2,4e-19 C impossivel), serie triboeletrica,
%  gerador de Van de Graaff e contagem de eletrons em fio de cobre.
%  Este banco explora o que restou: limite de eletrizacao pela rigidez
%  dieletrica, poder das pontas, ordem de grandeza do atrito, forca
%  eletrica x peso, conservacao de carga em processos nucleares,
%  gaiola/copo de Faraday e a serie geometrica dos contatos sucessivos.
%  Distribuicao mantida: 2 N3 + 8 N4 = 10
% =====================================================================

\subsubsection*{Banco complementar --- Cargas elétricas}

\begin{atencao}[Organização do banco]
Dois princípios sustentam o tópico: a carga é \textbf{conservada} (em sistema
isolado, a soma algébrica não muda) e é \textbf{quantizada} (todo valor é
múltiplo inteiro de $e=\pot{1.6}{-19}\,\si{\coulomb}$). O N4 explora as
consequências quantitativas menos óbvias --- há um teto para a carga que um
corpo pode reter, o atrito envolve uma fração ínfima dos átomos, e a
distribuição em condutores é governada pela curvatura.
\end{atencao}

\blocoaprofundamento

\begin{questao}[3]{}
Um sistema isolado contém $5$ prótons e $3$ elétrons livres. Após uma
reorganização interna, dois elétrons passam a orbitar dois dos prótons,
formando átomos neutros.
\begin{itensq}
\item Determine a carga total antes e depois.
\item Justifique o resultado.
\end{itensq}
\dados{$e=\pot{1.6}{-19}\,\si{\coulomb}$}
\resposta{$Q=+2e=\pot{3.2}{-19}\,\si{\coulomb}$ nos dois instantes}
\resolucao{(a) \textbf{Antes:} $Q=5(+e)+3(-e)=+2e=\pot{3.2}{-19}\,\si{\coulomb}$.\\
\textbf{Depois:} dois pares próton--elétron formam átomos neutros, restando
$3$ prótons livres e $1$ elétron livre:
\[
Q=3(+e)+1(-e)=+2e=\pot{3.2}{-19}\,\si{\coulomb}.
\]
(b) \textbf{Nada mudou} --- e não poderia mudar. A reorganização apenas
\emph{redistribuiu} as cargas dentro do sistema; nenhuma foi criada nem
destruída, e o sistema é isolado.\\
\textbf{Ponto conceitual:} um corpo \emph{neutro} não é um corpo \emph{sem
carga} --- ele tem tantas cargas positivas quanto negativas. Agrupar cargas em
átomos neutros esconde-as, mas não as elimina do balanço.}
\end{questao}

\begin{questao}[3]{}
Um bastão isolante adquire, por atrito, carga de $\SI{-5}{\nano\coulomb}$.
\begin{itensq}
\item Determine o número de elétrons transferidos.
\item Determine a massa correspondente e compare com a massa do bastão
      ($\SI{20}{\gram}$).
\end{itensq}
\dados{$e=\pot{1.6}{-19}\,\si{\coulomb}$;
$m_e=\pot{9.11}{-31}\,\si{\kilogram}$}
\resposta{$n=\pot{3.1}{10}$ elétrons; $\Delta m=\pot{2.8}{-20}\,\si{\kilogram}$}
\resolucao{(a)
\[
n=\frac{|Q|}{e}=\frac{\pot{5}{-9}}{\pot{1.6}{-19}}=\pot{3.125}{10}.
\]
Trinta bilhões de elétrons --- número enorme em termos absolutos.\\
(b)
\[
\Delta m=n\,m_e=\pot{3.125}{10}\times\pot{9.11}{-31}
=\pot{2.85}{-20}\,\si{\kilogram}.
\]
Comparando com os $\SI{20}{\gram}$ do bastão:
\[
\frac{\Delta m}{m}=\frac{\pot{2.85}{-20}}{\pot{2}{-2}}=\pot{1.4}{-18}.
\]
\textbf{Conclusão:} a eletrização por atrito \emph{não} altera a massa de forma
mensurável --- o desvio está $15$ ordens de grandeza abaixo da melhor balança
analítica. É por isso que a massa não serve como indicador de eletrização, e
por que o fenômeno passou séculos sendo descrito apenas por seus efeitos de
força.}
\end{questao}

\blocodesafio

\begin{questao}[4]{}
\textbf{(Série geométrica dos contatos.)} Uma esfera $A$ com carga $Q$ é
tocada sucessivamente por esferas idênticas \emph{neutras}
$B_1,B_2,\dots,B_n$, uma de cada vez, sendo cada uma retirada após o contato.
\begin{itensq}
\item Determine a carga de $A$ após $n$ contatos.
\item Determine a carga de cada esfera $B_k$.
\item Verifique a conservação da carga total.
\end{itensq}
\resposta{$Q_A=Q/2^{n}$; $Q_{B_k}=Q/2^{k}$; a soma reconstitui $Q$}
\resolucao{(a) Cada contato entre esferas \emph{idênticas} reparte igualmente a
carga presente. Como cada $B_k$ chega neutra, $A$ perde metade do que tem:
\[
Q\to\frac{Q}{2}\to\frac{Q}{4}\to\cdots
\;\Longrightarrow\;
Q_A=\frac{Q}{2^{n}}.
\]
(b) A esfera $B_k$ leva metade do que $A$ possuía imediatamente antes, ou seja
$\dfrac{1}{2}\cdot\dfrac{Q}{2^{k-1}}=\dfrac{Q}{2^{k}}$.\\
(c) \textbf{Conservação:}
\[
Q_A+\sum_{k=1}^{n}Q_{B_k}
=\frac{Q}{2^{n}}+Q\sum_{k=1}^{n}\frac{1}{2^{k}}
=\frac{Q}{2^{n}}+Q\left(1-\frac{1}{2^{n}}\right)=Q.\ \checkmark
\]
A soma geométrica $\sum 1/2^{k}$ tende a $1$, e o termo residual de $A$ fecha
exatamente a diferença.\\
\textbf{Leitura:} nenhuma carga se perde --- ela apenas se dilui. Com
$n=10$, $A$ retém menos de $0{,}1\%$ de $Q$, mas os $99{,}9\%$ restantes estão
distribuídos pelas dez esferas retiradas.\\
\textbf{Uso prático:} é assim que se \emph{descarrega} progressivamente um corpo
sem aterrá-lo, e também como se obtêm frações controladas de uma carga de
referência.}
\end{questao}

\begin{questao}[4]{}
\textbf{(Limite de eletrização.)} Uma esfera condutora de raio
$R=\SI{10}{\centi\meter}$ é carregada no ar, cuja rigidez dielétrica é
$E_{rup}=\SI{3e6}{\volt\per\meter}$.
\begin{itensq}
\item Determine a carga máxima que ela pode reter.
\item Determine o potencial correspondente.
\item Explique o que ocorre se tentarmos ultrapassar esse valor.
\end{itensq}
\dados{$\ke=\SI{9e9}{\newton\meter\squared\per\coulomb\squared}$}
\resposta{$Q_{\max}=\SI{3.33}{\micro\coulomb}$; $V=\SI{300}{\kilo\volt}$}
\resolucao{(a) O campo na superfície de uma esfera condutora vale
$E=\ke Q/R^{2}$. Impondo $E=E_{rup}$:
\[
Q_{\max}=\frac{E_{rup}R^{2}}{\ke}=\frac{\pot{3}{6}\times0{,}01}{\pot{9}{9}}
=\pot{3.33}{-6}\,\si{\coulomb}=\SI{3.33}{\micro\coulomb}.
\]
(b)
\[
V=\frac{\ke Q_{\max}}{R}=E_{rup}\,R=\pot{3}{6}(0{,}10)=\SI{300}{\kilo\volt}.
\]
\textbf{Relação notável:} $V_{\max}=E_{rup}R$ --- o potencial máximo é
\emph{proporcional ao raio}. Dobrar o raio dobra a tensão suportável, e é por
isso que geradores de Van de Graaff têm cúpulas grandes e lisas.\\
(c) \textbf{Acima do limite}, o ar ioniza e passa a conduzir: a carga escapa por
descarga corona (o brilho azulado e o chiado característicos) ou por faísca.
A esfera \emph{não consegue} reter mais carga --- ela se autolimita.\\
\textbf{Consequência de projeto:} para armazenar mais carga é preciso aumentar
$R$, usar meio de maior rigidez (óleo, $\mathrm{SF_6}$, vácuo) ou eliminar
irregularidades da superfície, onde o campo local é muito maior --- assunto da
questão seguinte.}
\end{questao}

\begin{questao}[4]{}
\textbf{(Poder das pontas.)} Duas esferas condutoras de raios
$R=\SI{10}{\centi\meter}$ e $r=\SI{1}{\centi\meter}$ são ligadas por um fio
longo e carregadas até $V=\SI{30}{\kilo\volt}$.
\begin{itensq}
\item Determine a razão entre as cargas e entre as densidades superficiais.
\item Determine o campo na superfície de cada uma.
\item Explique o efeito e cite duas aplicações.
\end{itensq}
\resposta{$Q_R/Q_r=10$; $\sigma_r/\sigma_R=10$;
$E_R=\SI{300}{\kilo\volt\per\meter}$ e $E_r=\SI{3}{\mega\volt\per\meter}$}
\resolucao{(a) Ligadas por fio, as duas ficam no \emph{mesmo potencial}:
\[
V=\frac{\ke Q}{R}\Rightarrow Q\propto R
\;\Longrightarrow\;
\frac{Q_R}{Q_r}=\frac{R}{r}=10.
\]
A densidade superficial é $\sigma=Q/(4\pi R^{2})\propto R/R^{2}=1/R$:
\[
\frac{\sigma_r}{\sigma_R}=\frac{R}{r}=10.
\]
\textbf{A esfera maior tem mais carga, mas a menor tem carga mais
concentrada.}\\
(b) $E=V/R$ na superfície de cada uma:
\[
E_R=\frac{\pot{30}{3}}{0{,}10}=\SI{300}{\kilo\volt\per\meter};
\qquad
E_r=\frac{\pot{30}{3}}{0{,}01}=\SI{3}{\mega\volt\per\meter}.
\]
(c) \textbf{O campo na esfera pequena atingiu exatamente a rigidez do ar.} A
ionização começa \emph{ali}, embora as duas estejam ao mesmo potencial ---
regiões de menor raio de curvatura concentram campo. É o \textbf{poder das
pontas}.\\
\textbf{Aplicações:}
\begin{itemize}
\item \textbf{Para-raios:} a ponta afiada ioniza o ar e estabelece o caminho
preferencial da descarga, protegendo a estrutura.
\item \textbf{Impressão e pintura eletrostática:} eletrodos pontiagudos ionizam
o ar para carregar tinta ou papel.
\end{itemize}
\textbf{E a contrapartida:} em equipamentos de alta tensão evita-se qualquer
ponta ou aresta viva --- daí os terminais arredondados e as blindagens
toroidais.}
\end{questao}

\begin{questao}[4]{}
\textbf{(Ordem de grandeza do atrito.)} Um bastão de vidro de $\SI{20}{\gram}$
adquire $\SI{5}{\nano\coulomb}$ por atrito. A massa molar média do vidro é
cerca de $\SI{60}{\gram\per\mole}$.
\begin{itensq}
\item Estime o número de átomos do bastão.
\item Determine a fração de átomos que perdeu um elétron.
\item Interprete o resultado.
\end{itensq}
\dados{$N_A=\pot{6.02}{23}\,\si{\per\mole}$}
\resposta{$\approx\pot{2}{23}$ átomos; fração de $\pot{1.6}{-13}$}
\resolucao{(a)
\[
N=\frac{m}{M}N_A=\frac{20}{60}\times\pot{6.02}{23}\approx\pot{2.0}{23}.
\]
(b) Já se sabe que $\pot{3.1}{10}$ elétrons foram transferidos:
\[
f=\frac{\pot{3.1}{10}}{\pot{2.0}{23}}=\pot{1.6}{-13},
\]
ou cerca de \textbf{um átomo em seis trilhões}.\\
(c) \textbf{Duas leituras que se complementam.}\\
\emph{Do ponto de vista microscópico}, a eletrização por atrito é um fenômeno
de superfície absolutamente marginal --- envolve uma fração desprezível da
matéria, e apenas a camada atômica mais externa.\\
\emph{Do ponto de vista macroscópico}, esses mesmos
$\pot{3.1}{10}$ elétrons produzem forças facilmente visíveis, capazes de
levantar pedaços de papel contra o peso.\\
\textbf{O contraste explica a intensidade da força elétrica:} ela é tão maior
que a gravitacional que um desequilíbrio de uma parte em $10^{13}$ já produz
efeitos observáveis. E explica também por que a matéria comum é neutra com
altíssima precisão --- qualquer desequilíbrio maior seria violentamente
desfeito.}
\end{questao}

\begin{questao}[4]{}
\textbf{(Força elétrica contra peso.)} Duas esferas idênticas de
$\SI{1}{\gram}$ estão a $\SI{5}{\centi\meter}$ uma da outra, cada uma com
carga $q$.
\begin{itensq}
\item Determine a força para $q=\SI{5}{\nano\coulomb}$ e para
      $q=\SI{50}{\nano\coulomb}$.
\item Determine a carga necessária para que a força iguale o peso.
\item Comente a viabilidade.
\end{itensq}
\dados{$g=\SI{9.8}{\meter\per\second\squared}$}
\resposta{$\SI{90}{\micro\newton}$ e $\SI{9}{\milli\newton}$;
$q\approx\SI{52}{\nano\coulomb}$}
\resolucao{(a)
\[
F=\frac{\ke q^{2}}{d^{2}}:\quad
q=\SI{5}{\nano\coulomb}\Rightarrow\SI{9.0e-5}{\newton};
\qquad
q=\SI{50}{\nano\coulomb}\Rightarrow\SI{9.0e-3}{\newton}.
\]
O peso vale $P=mg=\pot{1}{-3}(9{,}8)=\SI{9.8}{\milli\newton}$.\\
(b) Impondo $F=P$:
\[
q=d\sqrt{\frac{P}{\ke}}=0{,}05\sqrt{\frac{\pot{9.8}{-3}}{\pot{9}{9}}}
=0{,}05\times\pot{1.04}{-6}=\SI{52}{\nano\coulomb}.
\]
(c) \textbf{Perfeitamente viável.} Cinquenta nanocoulombs correspondem a
$\pot{3.3}{11}$ elétrons --- ordem de grandeza obtida rotineiramente por
atrito, como mostrou a questão anterior.\\
\textbf{É por isso que a eletrostática é observável a olho nu:} um simples
canudo atritado sustenta pedaços de papel contra a gravidade de todo o planeta.
\\
\textbf{Verificação do limite:} essas esferas suportariam essa carga? Para
$R=\SI{5}{\milli\meter}$, o campo na superfície seria
$\ke q/R^{2}=\SI{1.9}{\mega\volt\per\meter}$ --- abaixo da rigidez do ar, mas
já na mesma ordem. Cargas muito maiores escapariam por descarga.}
\end{questao}

\begin{questao}[4]{}
\textbf{(Conservação em processos nucleares.)} Verifique a conservação da carga
nos processos abaixo.
\begin{itensq}
\item Decaimento beta: $n\to p+e^{-}+\bar\nu_e$.
\item Aniquilação: $e^{-}+e^{+}\to2\gamma$.
\item Discuta o alcance do princípio.
\end{itensq}
\resposta{Conservada em ambos: $0=+e-e+0$ e $-e+e=0$}
\resolucao{(a) \textbf{Decaimento beta.} Antes: nêutron, carga $0$. Depois:
próton ($+e$), elétron ($-e$) e antineutrino ($0$):
\[
0=(+e)+(-e)+0=0.\ \checkmark
\]
Note que \emph{o número de partículas mudou} e que partículas foram
\emph{criadas} --- o elétron não estava dentro do nêutron. Ainda assim a carga
total se conserva.\\
(b) \textbf{Aniquilação.} Antes: $-e+(+e)=0$. Depois: dois fótons, ambos de
carga nula. \checkmark\\
Aqui a \emph{matéria} desaparece por completo, convertida em radiação --- e a
carga continua conservada, porque já era nula.\\
(c) \textbf{Alcance do princípio.} A conservação da carga é uma das leis mais
robustas da física:
\begin{itemize}
\item Vale em processos \textbf{químicos}, \textbf{nucleares} e de
\textbf{física de partículas}, sem exceção conhecida.
\item Vale mesmo quando partículas são criadas ou destruídas --- o que a
distingue da conservação do número de partículas, que \emph{não} é uma lei.
\item Experimentalmente, o decaimento do elétron (que violaria a conservação) foi
buscado e não observado; o limite atual para sua vida média é superior a
$10^{26}$ anos.
\end{itemize}
\textbf{Origem teórica:} a conservação da carga decorre de uma simetria
fundamental do eletromagnetismo, do mesmo modo que a conservação da energia
decorre da invariância no tempo. Ela não é um fato empírico isolado, mas
consequência da estrutura da teoria.}
\end{questao}

\begin{questao}[4]{}
\textbf{(Copo de Faraday.)} Um corpo carregado é introduzido, sem tocar as
paredes, no interior de um recipiente metálico fechado ligado a um eletrômetro.
\begin{itensq}
\item Explique o que o instrumento indica e por quê.
\item Explique por que a leitura não depende da posição do corpo dentro do
      recipiente.
\item Descreva o que ocorre se o corpo tocar a parede interna.
\end{itensq}
\resposta{O eletrômetro indica exatamente a carga do corpo, independentemente da
posição}
\resolucao{(a) \textbf{Indução total.} O corpo de carga $+Q$ induz $-Q$ na
superfície \emph{interna} do recipiente --- exatamente $-Q$, porque todas as
linhas de campo que saem do corpo terminam nela. Por conservação, sobra
$+Q$ na superfície \emph{externa}, e é essa carga que o eletrômetro mede.\\
\textbf{Resultado:} a leitura é $+Q$, o valor exato da carga introduzida.\\
(b) \textbf{Independência da posição.} O campo elétrico dentro do metal é nulo
em equilíbrio. Aplicando a lei de Gauss a uma superfície inteiramente contida no
metal, o fluxo é zero --- logo a carga interna total (corpo mais superfície
interna) é nula, \emph{qualquer que seja a posição do corpo}. A carga externa
fica sempre $+Q$, distribuída de acordo apenas com a geometria \emph{externa}
do recipiente.\\
\textbf{Consequência prática:} não é preciso centralizar o corpo nem controlar
sua orientação --- o que torna o copo de Faraday o método padrão de medição
absoluta de carga.\\
(c) \textbf{Se o corpo tocar a parede:} a carga $+Q$ neutraliza a
$-Q$ induzida, o corpo sai neutro e o recipiente fica com $+Q$ na superfície
externa. \textbf{A leitura do eletrômetro não muda} --- continua $+Q$.\\
Isso permite \emph{transferir} integralmente a carga de um corpo para o
recipiente, e é a base do método clássico de acumulação sucessiva de carga em
geradores eletrostáticos.}
\end{questao}

\begin{questao}[4]{}
\textbf{(Blindagem eletrostática.)} Uma malha metálica aterrada envolve um
instrumento sensível.
\begin{itensq}
\item Explique por que o campo externo não penetra.
\item Determine se a blindagem funciona também no sentido inverso.
\item Discuta a diferença entre malha aterrada e não aterrada.
\end{itensq}
\resposta{Blinda o interior contra campos externos; o aterramento é necessário
para blindar o exterior contra fontes internas}
\resolucao{(a) \textbf{Campo externo não penetra.} As cargas livres do condutor
se redistribuem até anular o campo em seu interior --- é a condição de
equilíbrio eletrostático. Qualquer campo externo induz uma distribuição
superficial cujo campo cancela exatamente o campo aplicado na região interna.\\
Isso vale \emph{mesmo sem aterramento}: a gaiola de Faraday isolada blinda o
interior. É o efeito que protege passageiros de um carro atingido por raio.\\
(b) \textbf{Sentido inverso: só com aterramento.} Uma carga $+Q$ colocada
\emph{dentro} da gaiola induz $-Q$ na face interna e $+Q$ na externa --- e essa
carga externa produz campo do lado de fora. \textbf{A blindagem falha.}\\
Com a malha \textbf{aterrada}, a carga $+Q$ da face externa escoa para a terra,
e o campo externo desaparece. Só então a blindagem funciona nos dois sentidos.\\
(c) \textbf{Resumo da diferença.}
\begin{center}
\begin{tabular}{lcc}
\hline
& não aterrada & aterrada\\
\hline
Protege o interior de campos externos & sim & sim\\
Impede que fontes internas afetem o exterior & não & sim\\
\hline
\end{tabular}
\end{center}
\textbf{Limitações práticas:} a malha só é eficaz para aberturas muito menores
que o comprimento de onda do sinal a blindar --- em eletrostática (campo
constante), qualquer malha fina basta; em radiofrequência, a exigência é
severa. E a impedância da conexão de terra importa: um fio longo de aterramento
tem indutância e degrada a blindagem em altas frequências.}
\end{questao}
